\scanpage{14}
\begin{center}- 9 -\end{center}
\hfill n° 57

\uline{Espace des sous-groupes fermés}, (n° 384)

C'est adopté, et ça ira en § 5 du chap. VIII d'Intégration. On vide l'espace uniforme des parties fermées et la prop. 3 (passent en exerc., avec appel d'exerc.).

Godement a fait une nouvelle rédaction en Congrès, qui a été relue, et améliorée quant aux virgules ; le rédacteur définitif a noté celles-ci.

\uline{Ce qui devrait être le chapitre des groupes de Lie} (n° 385 et 386).

Le n° 385 (``Distributions ponctuelles sur une variété''), sera incorporé dans le Fasc. Rés. des Variétés, après avoir été considérablement allégé. On n'y parlera pas du cas algébrique (ce qui simplifie en bas de la p. 3, le corps étant infini). Naturellement, ça doit être refondu avec les jets. On dira ``distributions à support fini''. Alléger les sorites (p. ex. celui des produits, n° 8 ; pourquoi pas les limites inductives et projectives ??).

L'énorme n° 386 (``Groupes formels et hyperalgèbres'') n'est pas du tout ce que Bourbaki voulait. D'une part le lien entre groupes analytiques et algèbres de Lie n'est fait que dans un mythique chapitre ultérieur. D'autre part on n'y démontre pas de résultats substantiels sur les groupes formels en caract. $p$. Ainsi Bourbaki vide toute la rédaction, et, comme Edith Piaf, ``on repart à zéro''. Cependant, le § 4 prévu par Cartier (correspondance entre groupes formels et algèbres de Lie en caract. 0, formule de Hausdorff), nettement plus substantiel que le reste, devra être rédigé pour usage ultérieur.

Là dessus Bourbaki s'aperçoit qu'il dispose d'une fort bonne rédaction de Koszul (``Variétés de groupes'', n° 340), sur laquelle on a intelligemment travaillé (voir Tribu n° 51). Ça doit facilement s'articuler avec le Fasc. Rés. Variétés. Bruhat devra rédiger sur cette base, sans chercher à faire d'originalité. Sinon le Livre de Lie est fichu !

\begin{center}-:-:-:-\end{center}
\begin{center}\uline{Séminaire BOURBAKI}\end{center}
Décembre 62 :

\diagram{\begin{tabular}{ll}
CARTIER & -- Fluctuations dans les suites de variables aléatoires indépendantes\\
\unsure{DEJHAN} & -- Papier Gårding à la S. M. F.\\
DIEUDONNÉ & -- Travail de Kostant (sur Amitsur -- \ldots{} etc)
\end{tabular}}
\mspage{13}
